\documentclass{ximera}

\input{../../../preamble.tex}


\definecolor{fvdnavy}{HTML}{071D43}
\definecolor{fvdcyan}{HTML}{20BFC6}
\definecolor{fvdcyanlight}{HTML}{EFFCFB}
\definecolor{fvdmagenta}{HTML}{F10D69}
\definecolor{fvdmagentalight}{HTML}{FFF1F7}
\definecolor{fvdtext}{HTML}{163044}





%=================================================
% Pedagogische omgevingen
% PDF: tcolorbox
% HTML: learning-box
%=================================================

\newenvironment{voorkennis}
{%
    \ifdefined\HCode
        \HCode{%
<section class="learning-box learning-box--prior">
<div class="learning-box__header">
<span class="learning-box__icon" aria-hidden="true"></span>
<span class="learning-box__title">Voorkennis</span>
</div>
<div class="learning-box__content">
        }%
        \begin{itemize}
    \else
       \begin{tcolorbox}[
    enhanced,
    breakable,
    colback=fvdcyanlight,
    colframe=fvdcyan,
    coltext=fvdtext,
    boxrule=0.5pt,
    leftrule=4pt,
    arc=4mm,
    outer arc=4mm,
    left=7mm,
    right=6mm,
    top=5mm,
    bottom=5mm,
    before skip=6mm,
    after skip=6mm,
    title=Voorkennis,
    coltitle=fvdcyan,
    fonttitle=\bfseries\Large,
    attach title to upper={\par\smallskip},
    boxed title style={
        empty,
        left=0mm,
        right=0mm,
        top=0mm,
        bottom=0mm
    },
    overlay={
        \draw[
            fvdcyan,
            line width=0.5pt
        ]
        ([xshift=42mm,yshift=-13mm]frame.north west)
        --
        ([xshift=-7mm,yshift=-13mm]frame.north east);
    }
]
        \begin{itemize}
    \fi
}
{%
    \end{itemize}
    \ifdefined\HCode
        \HCode{%
</div>
</section>
        }%
    \else
        \end{tcolorbox}
    \fi
}


\newenvironment{leerdoelen}
{%
    \ifdefined\HCode
        \HCode{%
<section class="learning-box learning-box--goals">
<div class="learning-box__header">
<span class="learning-box__icon" aria-hidden="true"></span>
<span class="learning-box__title">Leerdoelen</span>
</div>
<div class="learning-box__content">
        }%
        \begin{itemize}
    \else
\begin{tcolorbox}[
    enhanced,
    breakable,
    colback=fvdmagentalight,
    colframe=fvdmagenta,
    coltext=fvdtext,
    boxrule=0.5pt,
    leftrule=4pt,
    arc=4mm,
    outer arc=4mm,
    left=7mm,
    right=6mm,
    top=5mm,
    bottom=5mm,
    before skip=6mm,
    after skip=6mm,
    title=Leerdoelen,
    coltitle=fvdmagenta,
    fonttitle=\bfseries\Large,
    attach title to upper={\par\smallskip},
    boxed title style={
        empty,
        left=0mm,
        right=0mm,
        top=0mm,
        bottom=0mm
    },
    overlay={
        \draw[
            fvdmagenta,
            line width=0.5pt
        ]
        ([xshift=42mm,yshift=-13mm]frame.north west)
        --
        ([xshift=-7mm,yshift=-13mm]frame.north east);
    }
]
        \begin{itemize}
    \fi
}
{%
    \end{itemize}
    \ifdefined\HCode
        \HCode{%
</div>
</section>
        }%
    \else
        \end{tcolorbox}
    \fi
}


\addPrintStyle{../../..}

\title{Oefeningen — Functieonderzoek}
\author{Ingmar Herreman}

\begin{document}

% FVD_CHAPTER_DATA_START
% Automatisch gegenereerd uit index.tex.
% Niet handmatig aanpassen.
\fvdchapterdata{3}{Afleiden van samengestelde functie}{4}
% FVD_CHAPTER_DATA_END



\maketitle

\FVDbeginExercises

\begin{opmerking}
De oefeningen zijn opgebouwd van deelvaardigheden naar volledige analyse.

\textbf{\ESS} betekent dat de oefening noodzakelijk is om het
leerplandoel op het vereiste beheersingsniveau te bereiken. Ook meerdere
Verdiepingsopgaven zijn daarom essentieel: een functieonderzoek is meer
dan een afgeleide berekenen.
\end{opmerking}

% ============================================================
% BASIS 1
% ============================================================

\FVDsection{Basis — Rolle en Lagrange}

\begin{oefening}[Stelling van Rolle: voorwaarden controleren]{\oefB\ESS}
Onderzoek telkens of de stelling van Rolle toepasbaar is op het gegeven
interval. Als dat zo is, bepaal minstens één geschikte waarde \(c\).

\begin{question}
\[
f(x)=x^2-4x+3,\qquad [1,3].
\]
\begin{oplossing}
Veeltermfuncties zijn continu en afleidbaar op \(\mathbb{R}\).
Verder
\[
f(1)=0=f(3).
\]
Rolle is toepasbaar. Omdat
\[
f'(x)=2x-4,
\]
geldt
\[
f'(c)=0\Longleftrightarrow c=2.
\]
Dus
\[
\boxed{c=2}.
\]
\end{oplossing}
\end{question}

\begin{question}
\[
f(x)=\sin x,\qquad [0,\pi].
\]
\begin{oplossing}
\(\sin x\) is continu en afleidbaar en
\[
f(0)=f(\pi)=0.
\]
Dus Rolle is toepasbaar. Uit
\[
f'(x)=\cos x
\]
volgt
\[
\boxed{c=\frac{\pi}{2}}.
\]
\end{oplossing}
\end{question}

\begin{question}
\[
f(x)=\frac1x,\qquad [-1,1].
\]
\begin{oplossing}
Rolle is niet toepasbaar. De functie is niet eens gedefinieerd in
\(x=0\) en dus niet continu op \([-1,1]\).
\end{oplossing}
\end{question}

\begin{question}
\[
f(x)=\ln x,\qquad [1,e].
\]
\begin{oplossing}
De functie is continu op \([1,e]\) en afleidbaar op \((1,e)\), maar
\[
f(1)=0,\qquad f(e)=1.
\]
Omdat de eindwaarden niet gelijk zijn, is Rolle niet toepasbaar.
\end{oplossing}
\end{question}

\begin{question}
\[
f(x)=\cos x,\qquad [0,2\pi].
\]
\begin{oplossing}
De voorwaarden zijn vervuld en
\[
f(0)=1=f(2\pi).
\]
Verder
\[
f'(x)=-\sin x.
\]
Binnen \((0,2\pi)\) is
\[
\boxed{c=\pi}
\]
een geschikte waarde.
\end{oplossing}
\end{question}
\end{oefening}

\begin{oefening}[Middelwaardestelling van Lagrange]{\oefB\ESS}
Controleer de voorwaarden en bepaal alle gevraagde waarden \(c\).

\begin{question}
\[
f(x)=x^2,\qquad [1,3].
\]
\begin{oplossing}
De gemiddelde helling is
\[
\frac{f(3)-f(1)}{3-1}=\frac{9-1}{2}=4.
\]
Volgens Lagrange bestaat \(c\in(1,3)\) met
\[
f'(c)=4.
\]
Omdat \(f'(x)=2x\),
\[
2c=4\Longrightarrow\boxed{c=2}.
\]
\end{oplossing}
\end{question}

\begin{question}
\[
f(x)=\sqrt{x},\qquad [1,4].
\]
\begin{oplossing}
De functie is continu op \([1,4]\) en afleidbaar op \((1,4)\).
De gemiddelde helling is
\[
\frac{2-1}{4-1}=\frac13.
\]
Omdat
\[
f'(x)=\frac1{2\sqrt{x}},
\]
moet
\[
\frac1{2\sqrt{c}}=\frac13.
\]
Dus
\[
2\sqrt c=3
\Longrightarrow
\boxed{c=\frac94}.
\]
\end{oplossing}
\end{question}

\begin{question}
\[
f(x)=\ln x,\qquad [1,e].
\]
\begin{oplossing}
De gemiddelde helling is
\[
\frac{1-0}{e-1}=\frac1{e-1}.
\]
Omdat
\[
f'(x)=\frac1x,
\]
volgt
\[
\frac1c=\frac1{e-1}
\Longrightarrow
\boxed{c=e-1}.
\]
\end{oplossing}
\end{question}

\begin{question}
\[
f(x)=e^x,\qquad [0,1].
\]
\begin{oplossing}
De gemiddelde helling is
\[
e-1.
\]
Omdat \(f'(x)=e^x\),
\[
e^c=e-1
\Longrightarrow
\boxed{c=\ln(e-1)}.
\]
\end{oplossing}
\end{question}
\end{oefening}

% ============================================================
% BASIS 2
% ============================================================

\FVDsection{Basis — stijgen, dalen en extrema}

\begin{oefening}[Teken van \(f'\)]{\oefB\ESS}
Bepaal voor elke functie de intervallen waarop ze stijgt en daalt, en
bepaal de lokale extrema.

\begin{question}
\[
f(x)=\frac{x^3}{x^2-1}.
\]
\begin{oplossing}
\[
f'(x)=\frac{x^2(x^2-3)}{(x^2-1)^2}.
\]
De noemer is positief in het domein. Kritieke waarden:
\[
x=-\sqrt3,\ 0,\ \sqrt3,
\]
met \(x=\pm1\) buiten het domein.

Het teken van \(f'\) is positief voor \(|x|>\sqrt3\) en negatief voor
\(|x|<\sqrt3\), behalve aan de domeinonderbrekingen.

Dus er is een lokaal maximum bij \(x=-\sqrt3\) en een lokaal minimum
bij \(x=\sqrt3\). Bij \(x=0\) is geen extremum.
\end{oplossing}
\end{question}

\begin{question}
\[
f(x)=\frac{x^2-3}{x^2+1}.
\]
\begin{oplossing}
\[
f'(x)=\frac{2x(x^2+1)-2x(x^2-3)}{(x^2+1)^2}
=\frac{8x}{(x^2+1)^2}.
\]
De noemer is steeds positief. Dus \(f'<0\) voor \(x<0\) en \(f'>0\)
voor \(x>0\). De functie heeft een minimum bij
\[
x=0,\qquad f(0)=-3.
\]
Dus het minimum is
\[
\boxed{(0,-3)}.
\]
\end{oplossing}
\end{question}

\begin{question}
\[
f(x)=\frac{3}{x-3}.
\]
\begin{oplossing}
\[
f'(x)=-\frac{3}{(x-3)^2}<0
\]
voor elk \(x\neq3\). De functie is dus dalend op
\[
(-\infty,3)\quad\text{en}\quad(3,+\infty).
\]
Er zijn geen extrema.
\end{oplossing}
\end{question}

\begin{question}
\[
f(x)=x e^{-x}.
\]
\begin{oplossing}
\[
f'(x)=e^{-x}(1-x).
\]
Omdat \(e^{-x}>0\), is \(f'>0\) voor \(x<1\) en \(f'<0\) voor \(x>1\).
Dus er is een maximum bij
\[
\boxed{\left(1,\frac1e\right)}.
\]
\end{oplossing}
\end{question}

\begin{question}
\[
f(x)=\frac{\ln x}{x},\qquad x>0.
\]
\begin{oplossing}
\[
f'(x)=\frac{1-\ln x}{x^2}.
\]
De noemer is positief. Dus \(f'>0\) als \(\ln x<1\), dus voor \(x<e\),
en \(f'<0\) voor \(x>e\).

De functie heeft een maximum bij
\[
\boxed{\left(e,\frac1e\right)}.
\]
\end{oplossing}
\end{question}
\end{oefening}

\begin{oefening}[Horizontale raaklijn versus extremum]{\oefB\ESS}
Bepaal alle punten met horizontale raaklijn en beslis welke daarvan
extrema zijn.

\begin{question}
\[
f(x)=x^3.
\]
\begin{oplossing}
\[
f'(x)=3x^2.
\]
Dus alleen \(x=0\) geeft een horizontale raaklijn. Het teken van \(f'\)
is aan beide zijden positief, dus er is geen extremum. \((0,0)\) is een
stationair buigpunt.
\end{oplossing}
\end{question}

\begin{question}
\[
f(x)=x^4-4x^2.
\]
\begin{oplossing}
\[
f'(x)=4x^3-8x=4x(x^2-2).
\]
Stationaire waarden:
\[
x=-\sqrt2,\quad0,\quad\sqrt2.
\]
Tekenonderzoek van \(f'\) geeft:
bij \(x=-\sqrt2\) een minimum, bij \(x=0\) een maximum en bij
\(x=\sqrt2\) een minimum.

De functiewaarden zijn
\[
f(\pm\sqrt2)=-4,\qquad f(0)=0.
\]
\end{oplossing}
\end{question}

\begin{question}
\[
f(x)=\frac{x^3+5x^2+12x+8}{3x^2},\qquad x\neq0.
\]
\begin{oplossing}
Schrijf eerst
\[
f(x)=\frac13\left(x+5+\frac{12}{x}+\frac{8}{x^2}\right).
\]
Dan
\[
f'(x)=\frac13\left(1-\frac{12}{x^2}-\frac{16}{x^3}\right)
=\frac{x^3-12x-16}{3x^3}.
\]
Omdat
\[
x^3-12x-16=(x-4)(x+2)^2,
\]
zijn de stationaire waarden \(x=-2\) en \(x=4\).

Bij \(x=-2\) verandert het teken niet, dus geen extremum.
Bij \(x=4\) verandert \(f'\) van negatief naar positief, dus een lokaal
minimum.
\end{oplossing}
\end{question}
\end{oefening}

% ============================================================
% BASIS 3
% ============================================================

\FVDsection{Basis — tweede afgeleide en buigpunten}

\begin{oefening}[Kromming en buigpunten]{\oefB\ESS}
Onderzoek waar de grafiek van elke functie van kromming verandert en
bepaal de buigpunten.

\begin{question}
\[
f(x)=x^3-3x.
\]
\begin{oplossing}
\[
f''(x)=6x.
\]
Het teken verandert bij \(x=0\), dus
\[
\boxed{(0,0)}
\]
is een buigpunt.
\end{oplossing}
\end{question}

\begin{question}
\[
f(x)=\frac{x^2+4}{x}=x+\frac4x.
\]
\begin{oplossing}
\[
f'(x)=1-\frac4{x^2},
\qquad
f''(x)=\frac8{x^3}.
\]
De tweede afgeleide verandert van negatief naar positief bij \(x=0\),
maar \(x=0\) behoort niet tot het domein. Er is dus \textbf{geen}
buigpunt.
\end{oplossing}
\end{question}

\begin{question}
\[
f(x)=x e^{-x}.
\]
\begin{oplossing}
\[
f'(x)=e^{-x}(1-x),
\]
\[
f''(x)=e^{-x}(x-2).
\]
Omdat \(e^{-x}>0\), verandert \(f''\) van negatief naar positief bij
\(x=2\). Dus
\[
\boxed{\left(2,\frac2{e^2}\right)}
\]
is een buigpunt.
\end{oplossing}
\end{question}

\begin{question}
\[
f(x)=\ln x,\qquad x>0.
\]
\begin{oplossing}
\[
f''(x)=-\frac1{x^2}<0
\]
voor elk \(x>0\). De kromming verandert nergens, dus er zijn geen
buigpunten.
\end{oplossing}
\end{question}
\end{oefening}

% ============================================================
% VERDIEPING 1
% ============================================================

\FVDsection{Verdieping — volledig functieonderzoek}

\begin{oefening}[Volledig onderzoek van rationale functies]{\oefU\ESS}
Onderzoek volledig. Behandel waar relevant:

\begin{itemize}
  \item domein en continuïteit;
  \item nulpunten en teken;
  \item symmetrie;
  \item limieten en asymptoten;
  \item stijgen, dalen en extrema;
  \item kromming en buigpunten;
  \item synthese in een grafiekschets.
\end{itemize}

\begin{question}
\[
f(x)=-\frac1x.
\]
\begin{oplossing}
Domein:
\[
\mathbb{R}\setminus\{0\}.
\]
Geen nulpunten. De functie is oneven.

Verticale asymptoot:
\[
x=0.
\]
Horizontale asymptoot:
\[
y=0.
\]

\[
f'(x)=\frac1{x^2}>0
\]
op beide domeinintervallen, dus de functie stijgt op
\[
(-\infty,0)\quad\text{en}\quad(0,+\infty).
\]

\[
f''(x)=-\frac2{x^3}.
\]
De kromming verschilt links en rechts van \(0\), maar \(0\) ligt niet
in het domein, dus er is geen buigpunt.
\end{oplossing}
\end{question}

\begin{question}
\[
f(x)=\frac{x+4}{x+3}.
\]
\begin{oplossing}
Domein:
\[
\mathbb{R}\setminus\{-3\}.
\]
Nulpunt:
\[
x=-4.
\]
Verder
\[
f(x)=1+\frac1{x+3}.
\]
Dus verticale asymptoot \(x=-3\) en horizontale asymptoot \(y=1\).

\[
f'(x)=-\frac1{(x+3)^2}<0.
\]
Dus de functie is dalend op elk deel van het domein en heeft geen extrema.

\[
f''(x)=\frac2{(x+3)^3}.
\]
Het teken verandert over \(x=-3\), maar \(x=-3\) ligt niet in het domein.
Dus er is geen buigpunt.
\end{oplossing}
\end{question}

\begin{question}
\[
f(x)=\frac{x^2-4x+4}{x^2-1}.
\]
\begin{oplossing}
Domein:
\[
\mathbb{R}\setminus\{-1,1\}.
\]
De teller is \((x-2)^2\), dus \(x=2\) is een dubbel nulpunt.

Omdat teller en noemer dezelfde graad hebben, is
\[
y=1
\]
een horizontale asymptoot. Verder zijn
\[
x=-1,\qquad x=1
\]
verticale asymptoten.

De eerste afgeleide is
\[
f'(x)=\frac{4x^2-10x+4}{(x^2-1)^2}
=
\frac{2(2x-1)(x-2)}{(x^2-1)^2}.
\]
Het tekenonderzoek gebeurt met de scheidingswaarden
\[
-1,\quad \frac12,\quad1,\quad2.
\]
Daaruit volgen een lokaal maximum bij \(x=\frac12\) en een lokaal minimum
bij \(x=2\), rekening houdend met de domeinonderbrekingen.

Voor het volledige tweede-afgeleideonderzoek kan een CAS ter controle
worden gebruikt, maar de uiteindelijke conclusies moeten analytisch
worden verantwoord.
\end{oplossing}
\end{question}

\begin{question}
\[
f(x)=\frac{x^3}{x^2-1}.
\]
\begin{oplossing}
Dit is het volledige voorbeeld uit de theorie.

\[
\operatorname{dom}f=\mathbb{R}\setminus\{-1,1\},
\qquad
f(-x)=-f(x).
\]
Nulpunt: \(x=0\).

Verticale asymptoten:
\[
x=\pm1.
\]
Schuine asymptoot:
\[
y=x.
\]

\[
f'(x)=\frac{x^2(x^2-3)}{(x^2-1)^2}.
\]
Lokaal maximum bij
\[
x=-\sqrt3,
\]
lokaal minimum bij
\[
x=\sqrt3.
\]
Bij \(x=0\) geen extremum.

\[
f''(x)=\frac{2x(x^2+3)}{(x^2-1)^3}.
\]
De oorsprong is een stationair buigpunt.
\end{oplossing}
\end{question}
\end{oefening}

\begin{oefening}[Volledig onderzoek van andere functietypes]{\oefU\ESS}
Onderzoek het verloop en schets de grafiek.

\begin{question}
\[
f(x)=x\sqrt{x+2}.
\]
\begin{oplossing}
Domein:
\[
[-2,+\infty).
\]
Nulpunten:
\[
x=-2,\quad x=0.
\]

Voor \(x>-2\):
\[
f'(x)=\frac{3x+4}{2\sqrt{x+2}}.
\]
Dus \(f\) daalt op
\[
\left(-2,-\frac43\right)
\]
en stijgt op
\[
\left(-\frac43,+\infty\right).
\]
Het minimum ligt bij
\[
x=-\frac43.
\]
De functiewaarde is
\[
f\left(-\frac43\right)
=
-\frac43\sqrt{\frac23}.
\]
\end{oplossing}
\end{question}

\begin{question}
\[
f(x)=x e^{-x}.
\]
\begin{oplossing}
Domein \(\mathbb{R}\). Nulpunt \(x=0\).

\[
f'(x)=e^{-x}(1-x).
\]
Maximum bij
\[
\left(1,\frac1e\right).
\]

\[
f''(x)=e^{-x}(x-2).
\]
Buigpunt:
\[
\left(2,\frac2{e^2}\right).
\]

Verder
\[
\lim_{x\to+\infty}xe^{-x}=0,
\]
dus \(y=0\) is een horizontale asymptoot aan de rechterkant.
\end{oplossing}
\end{question}

\begin{question}
\[
f(x)=\frac{\ln x}{x}.
\]
\begin{oplossing}
Domein:
\[
(0,+\infty).
\]
Nulpunt:
\[
x=1.
\]

\[
f'(x)=\frac{1-\ln x}{x^2}.
\]
Maximum bij
\[
\left(e,\frac1e\right).
\]

De tweede afgeleide is
\[
f''(x)=\frac{2\ln x-3}{x^3}.
\]
Dus het buigpunt ligt bij
\[
2\ln x-3=0
\Longrightarrow
x=e^{3/2}.
\]
Daarbij hoort
\[
f(e^{3/2})=\frac{3}{2e^{3/2}}.
\]

Verder
\[
\lim_{x\to0^+}\frac{\ln x}{x}=-\infty
\]
en
\[
\lim_{x\to+\infty}\frac{\ln x}{x}=0.
\]
Dus \(x=0\) is een verticale asymptoot en \(y=0\) een horizontale
asymptoot voor \(x\to+\infty\).
\end{oplossing}
\end{question}
\end{oefening}

% ============================================================
% VERDIEPING 2
% ============================================================

\FVDsection{Verdieping — analyseren en redeneren}

\begin{oefening}[Van informatie naar grafiek]{\oefU\ESS}
Een functie \(f\) heeft domein \(\mathbb{R}\setminus\{2\}\) en voldoet
aan:

\[
\lim_{x\to2^-}f(x)=+\infty,
\qquad
\lim_{x\to2^+}f(x)=-\infty,
\]
\[
\lim_{x\to\pm\infty}f(x)=1,
\]
\[
f'(x)>0 \text{ voor }x<0,\qquad
f'(x)<0 \text{ voor }0<x<2,
\]
\[
f'(x)<0 \text{ voor }2<x<4,\qquad
f'(x)>0 \text{ voor }x>4.
\]

Verder is \(f(0)=3\) en \(f(4)=-2\).

\begin{question}
Bepaal de asymptoten.
\begin{oplossing}
Verticale asymptoot:
\[
\boxed{x=2}.
\]
Horizontale asymptoot:
\[
\boxed{y=1}.
\]
\end{oplossing}
\end{question}

\begin{question}
Bepaal de lokale extrema.
\begin{oplossing}
Bij \(x=0\) verandert \(f'\) van positief naar negatief:
\[
\boxed{(0,3)\text{ is een lokaal maximum}.}
\]
Bij \(x=4\) verandert \(f'\) van negatief naar positief:
\[
\boxed{(4,-2)\text{ is een lokaal minimum}.}
\]
\end{oplossing}
\end{question}

\begin{question}
Schets een mogelijke grafiek die aan alle gegevens voldoet.
\begin{oplossing}
De schets moet twee takken hebben, met verticale asymptoot \(x=2\),
horizontale asymptoot \(y=1\), een lokaal maximum in \((0,3)\) op de
linkertak en een lokaal minimum in \((4,-2)\) op de rechtertak.
\end{oplossing}
\end{question}
\end{oefening}

\begin{oefening}[Foutenanalyse]{\oefU\ESS}
Analyseer de redenering en corrigeer waar nodig.

\begin{question}
Een leerling schrijft:
``\(f'(2)=0\), dus \(f\) heeft een extremum bij \(x=2\).''
\begin{oplossing}
Onvoldoende. \(f'(2)=0\) toont alleen dat \(x=2\) stationair is.
Men moet bijvoorbeeld het teken van \(f'\) links en rechts van \(2\)
onderzoeken. Alleen bij een tekenverandering volgt een lokaal extremum.
\end{oplossing}
\end{question}

\begin{question}
Een leerling schrijft:
``\(f''(1)=0\), dus \((1,f(1))\) is een buigpunt.''
\begin{oplossing}
Onvoldoende. Men moet controleren of de kromming werkelijk verandert,
bijvoorbeeld via een tekenverandering van \(f''\).
\end{oplossing}
\end{question}

\begin{question}
Een leerling vindt bij een rationale functie \(f''(3)=0\), maar
\(3\notin\operatorname{dom}f\), en noemt \(x=3\) toch een buigpunt.
\begin{oplossing}
Dat is onmogelijk. Een buigpunt is een punt van de grafiek. Als
\(3\notin\operatorname{dom}f\), bestaat \((3,f(3))\) niet.
\end{oplossing}
\end{question}

\begin{question}
Een leerling gebruikt een grafiekprogramma en besluit dat een functie
geen extremum heeft omdat er in het venster \([-10,10]\times[-10,10]\)
geen extremum zichtbaar is.
\begin{oplossing}
Dit is geen geldige analytische conclusie. Het extremum kan buiten het
venster liggen of slecht zichtbaar zijn. Men moet het teken van \(f'\)
onderzoeken of het venster doelgericht aanpassen op basis van analytische
informatie.
\end{oplossing}
\end{question}
\end{oefening}

\begin{oefening}[Rolle en Lagrange als redeneerinstrument]{\oefU\ESS}
\begin{question}
Een functie \(f\) is continu op \([2,7]\), afleidbaar op \((2,7)\) en
voldoet aan \(f(2)=5\) en \(f(7)=5\). Wat kan je met zekerheid besluiten?

\begin{oplossing}
Volgens Rolle bestaat minstens één \(c\in(2,7)\) waarvoor
\[
\boxed{f'(c)=0}.
\]
Men kan niet besluiten dat dit punt uniek is of dat het een extremum is.
\end{oplossing}
\end{question}

\begin{question}
Een positie \(s(t)\) voldoet aan \(s(2)=15\,\mathrm m\) en
\(s(8)=51\,\mathrm m\). Veronderstel dat \(s\) continu is op \([2,8]\)
en afleidbaar op \((2,8)\). Toon aan dat er een ogenblik is waarop de
ogenblikkelijke snelheid \(6\,\mathrm{m/s}\) bedraagt.

\begin{oplossing}
De gemiddelde snelheid is
\[
\frac{51-15}{8-2}
=
\frac{36}{6}
=
6\,\mathrm{m/s}.
\]
Volgens Lagrange bestaat minstens één \(c\in(2,8)\) waarvoor
\[
s'(c)=6\,\mathrm{m/s}.
\]
\end{oplossing}
\end{question}

\begin{question}
Waarom kan dezelfde conclusie niet automatisch worden getrokken als de
plaatsfunctie op één tijdstip tussen \(2\) en \(8\) een sprong maakt?

\begin{oplossing}
Dan is de continuïteitsvoorwaarde op \([2,8]\) geschonden. De
middelwaardestelling van Lagrange is dan niet toepasbaar.
\end{oplossing}
\end{question}
\end{oefening}

% ============================================================
% EXTREMUMPROBLEMEN
% ============================================================

\FVDsection{Verdieping — extremumproblemen}

\begin{oefening}[Van context naar doelfunctie]{\oefU\ESS}
Los de problemen volledig op. Noteer expliciet:

\begin{itemize}
  \item de gekozen variabele;
  \item de doelfunctie;
  \item het betekenisvolle domein;
  \item de extremumberekening;
  \item de interpretatie met eenheden.
\end{itemize}

\begin{question}
Een rechthoek heeft omtrek \(60\,\mathrm m\). Bepaal de afmetingen
waarvoor de oppervlakte maximaal is.

\begin{oplossing}
Neem lengte \(x\) en breedte \(y\). Uit
\[
2x+2y=60
\]
volgt
\[
y=30-x.
\]
Dus
\[
A(x)=x(30-x)=30x-x^2,
\qquad 0<x<30.
\]
\[
A'(x)=30-2x.
\]
Daarom
\[
A'(x)=0\Longleftrightarrow x=15.
\]
Omdat \(A''(x)=-2<0\), is dit een maximum.

Dus
\[
\boxed{15\,\mathrm m\times15\,\mathrm m}
\]
met maximale oppervlakte
\[
\boxed{225\,\mathrm m^2}.
\]
\end{oplossing}
\end{question}

\begin{question}
Een rechthoek wordt ingeschreven onder de parabool
\[
y=12-x^2
\]
met de basis op de \(x\)-as en symmetrisch rond de \(y\)-as. Bepaal de
afmetingen van de rechthoek met maximale oppervlakte.

\begin{oplossing}
Neem \(x>0\) als halve breedte. Dan is de volledige breedte \(2x\) en
de hoogte
\[
12-x^2.
\]
De doelfunctie is
\[
A(x)=2x(12-x^2)=24x-2x^3.
\]
Uit de context:
\[
0<x<\sqrt{12}.
\]
\[
A'(x)=24-6x^2.
\]
Dus
\[
x^2=4\Longrightarrow x=2.
\]
De breedte is \(4\) en de hoogte
\[
12-4=8.
\]
Dus de optimale afmetingen zijn
\[
\boxed{4\times8}
\]
en de maximale oppervlakte is
\[
\boxed{32}.
\]
\end{oplossing}
\end{question}

\begin{question}
Uit een vierkant karton van \(24\,\mathrm{cm}\) bij
\(24\,\mathrm{cm}\) worden in elke hoek vierkantjes met zijde \(x\)
geknipt. De randen worden omhoog gevouwen tot een open doos. Voor welke
\(x\) is het volume maximaal?

\begin{oplossing}
\[
V(x)=x(24-2x)^2,
\qquad 0<x<12.
\]
\[
V'(x)=(24-2x)(24-6x).
\]
Kandidaten zijn \(x=12\) en \(x=4\), maar \(x=12\) is een randpunt waar
het volume nul wordt. Binnen het domein geeft
\[
\boxed{x=4\,\mathrm{cm}}
\]
het maximum.

De doos is dan
\[
16\,\mathrm{cm}\times16\,\mathrm{cm}\times4\,\mathrm{cm}
\]
en
\[
\boxed{V_{\max}=1024\,\mathrm{cm}^3}.
\]
\end{oplossing}
\end{question}
\end{oefening}

\begin{oefening}[STEM — maximale vermogensoverdracht]{\oefU\ESS}
Een bron met elektromotorische spanning \(E>0\) en inwendige weerstand
\(r>0\) levert aan een belasting \(R>0\) het vermogen

\[
P(R)=\frac{E^2R}{(R+r)^2}.
\]

\begin{question}
Bepaal \(P'(R)\).

\begin{oplossing}
\[
P'(R)
=
E^2\frac{(R+r)^2-2R(R+r)}{(R+r)^4}
=
E^2\frac{r-R}{(R+r)^3}.
\]
\end{oplossing}
\end{question}

\begin{question}
Voor welke \(R\) is het vermogen maximaal?

\begin{oplossing}
Voor \(R>0\) zijn \(E^2\) en \((R+r)^3\) positief. Het teken van \(P'\)
wordt dus bepaald door \(r-R\).

Daarom:
\[
P'(R)>0\quad\text{voor }R<r,
\]
\[
P'(R)<0\quad\text{voor }R>r.
\]
Dus
\[
\boxed{R=r}
\]
geeft het maximale vermogen.
\end{oplossing}
\end{question}

\begin{question}
Bepaal het maximale vermogen.

\begin{oplossing}
\[
P(r)
=
\frac{E^2r}{(2r)^2}
=
\boxed{\frac{E^2}{4r}}.
\]
\end{oplossing}
\end{question}
\end{oefening}

% ============================================================
% GEVORDERD
% ============================================================

\FVDsection{Gevorderd — parameters, modelleren en synthese}

\begin{oefening}[Parameteronderzoek]{\oefG}
Gegeven is

\[
f_a(x)=\frac{x-2}{1-ax}.
\]

Onderzoek voor welke waarden van \(a\) de functie op elk interval van
haar domein strikt stijgend is.

\begin{oplossing}
Met de quotiëntregel:
\[
f_a'(x)
=
\frac{1(1-ax)-(x-2)(-a)}{(1-ax)^2}
=
\frac{1-2a}{(1-ax)^2}.
\]

Op het domein is de noemer positief. De functie is dus strikt stijgend
als
\[
1-2a>0.
\]
Daaruit volgt
\[
\boxed{a<\frac12}.
\]

Voor \(a=\frac12\) is de afgeleide nul op het domein en is de functie
constant op elk domeininterval; voor \(a>\frac12\) is ze dalend.
\end{oplossing}
\end{oefening}

\begin{oefening}[Parameter en extremum]{\oefG}
Bepaal \(a\) en \(b\) zodat

\[
f(x)=\frac{ax}{x^2+b}
\]

bij \(x=\frac52\) een lokale extreme waarde \(1\) bereikt.

\begin{oplossing}
De voorwaarden zijn
\[
f\left(\frac52\right)=1
\]
en
\[
f'\left(\frac52\right)=0.
\]

Eerst
\[
f'(x)
=
a\frac{(x^2+b)-2x^2}{(x^2+b)^2}
=
a\frac{b-x^2}{(x^2+b)^2}.
\]

Voor een niet-triviale oplossing \(a\neq0\) geeft
\[
b-\left(\frac52\right)^2=0,
\]
dus
\[
\boxed{b=\frac{25}{4}}.
\]

Dan
\[
1
=
f\left(\frac52\right)
=
\frac{a\cdot\frac52}{\frac{25}{4}+\frac{25}{4}}
=
\frac{a\cdot\frac52}{\frac{25}{2}}
=
\frac a5.
\]
Dus
\[
\boxed{a=5}.
\]

Omdat het teken van \(f'\) bij \(x=\frac52\) van positief naar negatief
verandert, is het inderdaad een lokaal maximum.
\end{oplossing}
\end{oefening}

\begin{oefening}[Ontwerp een functieonderzoek]{\oefG}
Kies zelf één functie uit elk van de volgende twee groepen:

\begin{itemize}
  \item een rationale functie met minstens één verticale asymptoot;
  \item een exponentiële of logaritmische functie met minstens één extremum of buigpunt.
\end{itemize}

Voer voor beide een volledig functieonderzoek uit. Gebruik daarna een
grafiekprogramma uitsluitend om je resultaat te controleren.

Beschrijf ten minste één verschil in aanpak tussen beide functietypes.

\begin{oplossing}
Dit is een open onderzoeksopdracht. Een volledige oplossing moet minstens
bevatten:

\begin{itemize}
  \item correcte domeinanalyse;
  \item relevante limieten en asymptoten;
  \item correcte eerste en tweede afgeleide;
  \item tekenanalyse van \(f'\) en \(f''\);
  \item extrema en buigpunten;
  \item een grafiekschets;
  \item een vergelijking tussen de aanpak voor beide functietypes;
  \item een expliciete controle met ICT zonder de analytische redenering
  door de software te vervangen.
\end{itemize}
\end{oplossing}
\end{oefening}

\begin{oefening}[Modelleerprobleem zonder vooraf gegeven functie]{\oefG}
Een cilindervormig blik moet een vast volume \(V>0\) hebben.
Veronderstel dat de materiaalkost evenredig is met de totale oppervlakte
van een gesloten cilinder.

Bepaal de verhouding tussen hoogte \(h\) en straal \(r\) waarvoor zo
weinig mogelijk materiaal nodig is.

\begin{oplossing}
De volumebeperking is
\[
V=\pi r^2h,
\]
dus
\[
h=\frac{V}{\pi r^2}.
\]

De oppervlakte is
\[
A=2\pi r^2+2\pi rh.
\]

Invullen geeft
\[
A(r)=2\pi r^2+\frac{2V}{r},
\qquad r>0.
\]

De afgeleide is
\[
A'(r)=4\pi r-\frac{2V}{r^2}.
\]

Voor een extremum:
\[
4\pi r-\frac{2V}{r^2}=0
\]
\[
4\pi r^3=2V
\]
\[
2\pi r^3=V.
\]

Omdat
\[
V=\pi r^2h,
\]
volgt
\[
\pi r^2h=2\pi r^3.
\]

Dus
\[
\boxed{h=2r}.
\]

De optimale hoogte is dus gelijk aan de diameter van de cilinder.
Verder
\[
A''(r)=4\pi+\frac{4V}{r^3}>0,
\]
dus het gevonden punt geeft een minimum.
\end{oplossing}
\end{oefening}

\begin{opmerking}[Minimale beheersing]
Voor het officiële beheersingsniveau moeten minstens alle oefeningen met
\ESS beheerst zijn.

Dat omvat:

\begin{itemize}
  \item Rolle en Lagrange toepassen na controle van de voorwaarden;
  \item tekenonderzoek van \(f'\) en \(f''\);
  \item extrema en buigpunten correct classificeren;
  \item volledige functieonderzoeken uitvoeren voor meer dan alleen
  veeltermfuncties;
  \item informatie synthetiseren tot een grafiek;
  \item fouten in redeneringen analyseren;
  \item en extremumproblemen vanuit een context mathematiseren en oplossen.
\end{itemize}

De \oefG-opgaven gaan verder in zelfstandigheid, parametrisering,
modellering en voorbereiding op verdere studies.
\end{opmerking}

\FVDendExercises

\end{document}

